% ==============================================================================
% BỘ ĐỀ THAM KHẢO TUYỂN SINH 10 NĂM HỌC 2026-2027 SỞ GD&ĐT TP.HCM
% ĐỢT 19: CỤM 13 (ĐỀ 1 & ĐỀ 2 - ĐỀ SỐ 37 VÀ ĐỀ SỐ 38)
% Gồm 14 bài toán tự luận chuẩn cấu trúc tuyển sinh 10
% ==============================================================================

% ==============================================================================
% CỤM 13 - ĐỀ THAM KHẢO 1 (ĐỀ SỐ 37)
% ==============================================================================

\begin{bt}%[9D6H18-2]
% nguon: Đề tham khảo Tuyển sinh 10 năm học 2026-2027 Sở GD&ĐT TP.HCM - Cụm 13 Đề 1 (Đề số 37)
Cho hàm số $y = \dfrac{1}{2}x^2$ có đồ thị là parabol $(P)$.
\begin{enumerate}[a)]
	\item Vẽ đồ thị $(P)$ trên mặt phẳng toạ độ $Oxy$.
	\item Tìm toạ độ các điểm thuộc đồ thị $(P)$ sao cho hoành độ và tung độ có tổng bằng $12$.
\end{enumerate}
\loigiai{
\begin{enumerate}[a)]
	\item Bảng giá trị tương ứng của $x$ và $y$:
	\begin{center}
		\begin{tabular}{|c|c|c|c|c|c|}
			\hline
			$x$ & $-4$ & $-2$ & $0$ & $2$ & $4$ \\ \hline
			$y = \dfrac{1}{2}x^2$ & $8$ & $2$ & $0$ & $2$ & $8$ \\ \hline
		\end{tabular}
	\end{center}
	Đồ thị $(P)$ là một đường parabol đi qua các điểm $(-4; 8)$, $(-2; 2)$, $(0; 0)$, $(2; 2)$, $(4; 8)$, nhận trục tung $Oy$ làm trục đối xứng và có bề lõm quay lên trên.
	\begin{center}
	\begin{tikzpicture}[>=stealth,scale=0.7]
		\draw[->] (-5,0) -- (5,0) node[below] {$x$};
		\draw[->] (0,-1) -- (0,9) node[left] {$y$};
		\node[below left] at (0,0) {$O$};
		\draw[dashed] (-4,0) node[below] {$-4$} -- (-4,8) -- (0,8) node[left] {$8$} -- (4,8) -- (4,0) node[below] {$4$};
		\draw[dashed] (-2,0) node[below] {$-2$} -- (-2,2) -- (0,2) node[left] {$2$} -- (2,2) -- (2,0) node[below] {$2$};
		\draw[thick,smooth,samples=100,domain=-4.2:4.2] plot(\x,{0.5*(\x)^2});
		\fill (-4,8) circle (1.5pt) (-2,2) circle (1.5pt) (0,0) circle (1.5pt) (2,2) circle (1.5pt) (4,8) circle (1.5pt);
		\node[right] at (2.5,5) {$(P)\colon y = \dfrac{1}{2}x^2$};
	\end{tikzpicture}
	\end{center}
	\item Gọi $(x; y)$ là toạ độ điểm thuộc đồ thị $(P)$.
	Vì tổng của hoành độ và tung độ bằng $12$ nên:
	$$x + y = 12 \implies y = 12 - x.$$
	Thay vào phương trình của $(P)$, ta có:
	$$12 - x = \dfrac{1}{2}x^2 \iff x^2 + 2x - 24 = 0.$$
	Biệt thức $\Delta' = 1^2 - 1 \cdot (-24) = 25 > 0 \implies \sqrt{\Delta'} = 5$.
	Phương trình có hai nghiệm:
	$$x_1 = -1 + 5 = 4 \implies y_1 = 12 - 4 = 8;$$
	$$x_2 = -1 - 5 = -6 \implies y_2 = 12 - (-6) = 18.$$
	Vậy toạ độ các điểm cần tìm là $(4; 8)$ và $(-6; 18)$.
\end{enumerate}
}
\end{bt}

\begin{bt}%[9D6H20-2]
% nguon: Đề tham khảo Tuyển sinh 10 năm học 2026-2027 Sở GD&ĐT TP.HCM - Cụm 13 Đề 1 (Đề số 37)
Cho phương trình $3x^2 - 11x - 15 = 0$.
\begin{enumerate}[a)]
	\item Chứng minh phương trình luôn có hai nghiệm phân biệt $x_1, x_2$.
	\item Không giải phương trình, hãy tính giá trị của biểu thức $A = \dfrac{3x_1}{x_2} + \dfrac{3x_2}{x_1}$.
\end{enumerate}
\loigiai{
\begin{enumerate}[a)]
	\item Phương trình $3x^2 - 11x - 15 = 0$ có $a = 3$, $b = -11$, $c = -15$.
	Ta có:
	$$\Delta = b^2 - 4ac = (-11)^2 - 4 \cdot 3 \cdot (-15) = 121 + 180 = 301 > 0.$$
	Vì $\Delta > 0$ nên phương trình luôn có hai nghiệm phân biệt $x_1, x_2$.
	\item Áp dụng định lý Viète, ta có:
	$$\begin{cases} x_1 + x_2 = -\dfrac{b}{a} = \dfrac{11}{3} \\ x_1 x_2 = \dfrac{c}{a} = \dfrac{-15}{3} = -5. \end{cases}$$
	Ta có:
	$$x_1^2 + x_2^2 = (x_1 + x_2)^2 - 2x_1 x_2 = \left(\dfrac{11}{3}\right)^2 - 2(-5) = \dfrac{121}{9} + 10 = \dfrac{211}{9}.$$
	Khi đó:
	$$A = \dfrac{3x_1}{x_2} + \dfrac{3x_2}{x_1} = \dfrac{3(x_1^2 + x_2^2)}{x_1 x_2} = \dfrac{3 \cdot \dfrac{211}{9}}{-5} = \dfrac{\dfrac{211}{3}}{-5} = -\dfrac{211}{15}.$$
	Vậy $A = -\dfrac{211}{15}$.
\end{enumerate}
}
\end{bt}

\begin{bt}%[9D8H26-3]
% nguon: Đề tham khảo Tuyển sinh 10 năm học 2026-2027 Sở GD&ĐT TP.HCM - Cụm 13 Đề 1 (Đề số 37)
Một đĩa tròn bằng bìa cứng được chia làm $10$ phần bằng nhau và ghi các số $1; 2; 3; 4; 5; 6; 7; 8; 9; 10$. Chiếc kim được gắn cố định vào trục quay ở tâm đĩa. Người ta quay đĩa tròn một lần.
\begin{enumerate}[a)]
	\item Mô tả không gian mẫu của phép thử và cho biết không gian mẫu có bao nhiêu phần tử.
	\item Tính xác suất của mỗi biến cố sau:
	\begin{itemize}
		\item $A$: ``Mũi tên chỉ vào hình quạt ghi số là hợp số''.
		\item $B$: ``Mũi tên chỉ vào hình quạt ghi số là số chính phương''.
		\item $C$: ``Mũi tên chỉ vào hình quạt ghi số chia hết cho $5$''.
	\end{itemize}
\end{enumerate}
\loigiai{
\begin{enumerate}[a)]
	\item Không gian mẫu của phép thử là tập hợp các số ghi trên đĩa tròn:
	$$\Omega = \{1; 2; 3; 4; 5; 6; 7; 8; 9; 10\}.$$
	Số phần tử của không gian mẫu là $n(\Omega) = 10$.
	\item Tính xác suất của các biến cố:
	\begin{itemize}
		\item Xét biến cố $A$: ``Mũi tên chỉ vào hình quạt ghi số là hợp số''.
		Trong các số từ $1$ đến $10$, các hợp số là $4; 6; 8; 9; 10$ (số $1$ không phải là số nguyên tố cũng không phải là hợp số).
		Số kết quả thuận lợi cho biến cố $A$ là $n(A) = 5$.
		Xác suất của biến cố $A$ là:
		$$P(A) = \dfrac{n(A)}{n(\Omega)} = \dfrac{5}{10} = \dfrac{1}{2}.$$
		\item Xét biến cố $B$: ``Mũi tên chỉ vào hình quạt ghi số là số chính phương''.
		Các số chính phương trong không gian mẫu là $1; 4; 9$.
		Số kết quả thuận lợi cho biến cố $B$ là $n(B) = 3$.
		Xác suất của biến cố $B$ là:
		$$P(B) = \dfrac{n(B)}{n(\Omega)} = \dfrac{3}{10}.$$
		\item Xét biến cố $C$: ``Mũi tên chỉ vào hình quạt ghi số chia hết cho $5$''.
		Các số chia hết cho $5$ trong không gian mẫu là $5; 10$.
		Số kết quả thuận lợi cho biến cố $C$ là $n(C) = 2$.
		Xác suất của biến cố $C$ là:
		$$P(C) = \dfrac{n(C)}{n(\Omega)} = \dfrac{2}{10} = \dfrac{1}{5}.$$
	\end{itemize}
\end{enumerate}
}
\end{bt}

\begin{bt}%[9D6V21-3]
% nguon: Đề tham khảo Tuyển sinh 10 năm học 2026-2027 Sở GD&ĐT TP.HCM - Cụm 13 Đề 1 (Đề số 37)
Một công viên hình chữ nhật có chiều rộng là $x$ (m), chiều dài gấp rưỡi chiều rộng. Ở giữa công viên, người ta xây một đài phun nước hình tròn có đường kính bằng chiều rộng của công viên.
\begin{enumerate}[a)]
	\item Hãy viết biểu thức $S$ biểu diễn diện tích phần còn lại của công viên (phần không có đài phun nước) theo $x$ và $\pi$.
	\item Biết diện tích phần còn lại của công viên là $471{,}5\text{ m}^2$. Tính kích thước của công viên (lấy $\pi \approx 3{,}14$).
\end{enumerate}
\loigiai{
\begin{enumerate}[a)]
	\item Chiều rộng của công viên là $x$ (m) ($x > 0$).
	Chiều dài của công viên là $1{,}5x$ (m).
	Diện tích toàn bộ công viên là:
	$$S_1 = x \cdot 1{,}5x = 1{,}5x^2\text{ (m}^2\text{)}.$$
	Đài phun nước hình tròn có đường kính là $x$ (m) nên bán kính là $R = \dfrac{x}{2}$ (m).
	Diện tích đài phun nước là:
	$$S_2 = \pi R^2 = \pi \left(\dfrac{x}{2}\right)^2 = \dfrac{\pi}{4}x^2\text{ (m}^2\text{)}.$$
	Biểu thức biểu diễn diện tích phần còn lại của công viên là:
	$$S = S_1 - S_2 = 1{,}5x^2 - \dfrac{\pi}{4}x^2 = \left(1{,}5 - \dfrac{\pi}{4}\right)x^2\text{ (m}^2\text{)}.$$
	\item Theo giả thiết, diện tích phần còn lại là $471{,}5\text{ m}^2$ với $\pi \approx 3{,}14$, ta có phương trình:
	$$\left(1{,}5 - \dfrac{3{,}14}{4}\right)x^2 = 471{,}5$$
	$$\iff (1{,}5 - 0{,}785)x^2 = 471{,}5$$
	$$\iff 0{,}715x^2 = 471{,}5$$
	$$\iff x^2 = \dfrac{471{,}5}{0{,}715} = \dfrac{94300}{143}.$$
	Vì $x > 0$ nên $x = \sqrt{\dfrac{94300}{143}}\text{ (m)} \approx 25{,}68\text{ (m)}$.
	Chiều dài của công viên là:
	$$1{,}5x = 1{,}5 \cdot \sqrt{\dfrac{94300}{143}} = \dfrac{3}{2}\sqrt{\dfrac{94300}{143}}\text{ (m)} \approx 38{,}53\text{ (m)}.$$
	Vậy công viên có chiều rộng là $\sqrt{\dfrac{94300}{143}}\text{ m}$ và chiều dài là $\dfrac{3}{2}\sqrt{\dfrac{94300}{143}}\text{ m}$.
\end{enumerate}
}
\end{bt}

\begin{bt}%[9H0V31-3]
% nguon: Đề tham khảo Tuyển sinh 10 năm học 2026-2027 Sở GD&ĐT TP.HCM - Cụm 13 Đề 1 (Đề số 37)
Một bình chứa hình trụ có đường kính đáy $1\text{ dm}$, chiều cao $0{,}8\text{ dm}$. Bên trong bình có chứa sẵn một viên bi hình cầu có bán kính $3\text{ cm}$. Hỏi phải đổ vào bình bao nhiêu lít nước để nước đầy bình? (Làm tròn kết quả đến chữ số thập phân thứ nhất).
\loigiai{
Đổi đơn vị: $3\text{ cm} = 0{,}3\text{ dm}$.
Bình hình trụ có bán kính đáy là:
$$r = 1 : 2 = 0{,}5\text{ (dm)}.$$
Thể tích của bình hình trụ là:
$$V_{\text{trụ}} = \pi r^2 h = \pi \cdot (0{,}5)^2 \cdot 0{,}8 = 0{,}2\pi\text{ (dm}^3\text{)}.$$
Thể tích của viên bi hình cầu là:
$$V_{\text{cầu}} = \dfrac{4}{3}\pi R^3 = \dfrac{4}{3}\pi \cdot (0{,}3)^3 = 0{,}036\pi\text{ (dm}^3\text{)}.$$
Thể tích nước cần đổ vào bình để đầy bình là:
$$V_{\text{nước}} = V_{\text{trụ}} - V_{\text{cầu}} = 0{,}2\pi - 0{,}036\pi = 0{,}164\pi\text{ (dm}^3\text{)}.$$
Vì $1\text{ dm}^3 = 1\text{ lít}$ nên thể tích nước cần đổ vào là:
$$V_{\text{nước}} = 0{,}164\pi\text{ (lít)} \approx 0{,}164 \cdot 3{,}1416 \approx 0{,}515 \approx 0{,}5\text{ (lít)}.$$
Vậy cần đổ vào bình khoảng $0{,}5$ lít nước để đầy bình.
}
\end{bt}

\begin{bt}%[9D1H3-4]
% nguon: Đề tham khảo Tuyển sinh 10 năm học 2026-2027 Sở GD&ĐT TP.HCM - Cụm 13 Đề 1 (Đề số 37)
Bác Tư đến siêu thị mua một cái quạt máy và một ấm đun siêu tốc với tổng số tiền theo giá niêm yết là $630$ nghìn đồng. Tuy nhiên, trong tuần lễ tri ân khách hàng nên siêu thị đã giảm giá quạt máy $15\%$ và giảm giá ấm đun siêu tốc $12\%$ so với giá niêm yết của từng sản phẩm. Do đó Bác Tư chỉ phải trả $543$ nghìn đồng khi mua hai sản phẩm trên. Hỏi giá niêm yết ban đầu của mỗi sản phẩm là bao nhiêu?
\loigiai{
Gọi $x$ (nghìn đồng) là giá niêm yết của quạt máy và $y$ (nghìn đồng) là giá niêm yết của ấm đun siêu tốc ($0 < x, y < 630$).
Vì tổng số tiền theo giá niêm yết là $630$ nghìn đồng nên ta có phương trình:
$$x + y = 630. \quad (1)$$
Số tiền Bác Tư phải trả cho chiếc quạt máy sau khi giảm giá $15\%$ là:
$$(100\% - 15\%)x = 0{,}85x\text{ (nghìn đồng)}.$$
Số tiền Bác Tư phải trả cho chiếc ấm đun siêu tốc sau khi giảm giá $12\%$ là:
$$(100\% - 12\%)y = 0{,}88y\text{ (nghìn đồng)}.$$
Vì tổng số tiền thực tế Bác Tư phải trả là $543$ nghìn đồng nên ta có phương trình:
$$0{,}85x + 0{,}88y = 543. \quad (2)$$
Từ $(1)$ và $(2)$ ta có hệ phương trình:
$$\begin{cases} x + y = 630 \\ 0{,}85x + 0{,}88y = 543 \end{cases} \iff \begin{cases} 0{,}85x + 0{,}85y = 535{,}5 \\ 0{,}85x + 0{,}88y = 543 \end{cases} \iff \begin{cases} 0{,}03y = 7{,}5 \\ x + y = 630 \end{cases} \iff \begin{cases} y = 250 \\ x = 380 \end{cases}\text{ (thoả mãn)}.$$
Vậy giá niêm yết của quạt máy là $380\,000$ đồng và của ấm đun siêu tốc là $250\,000$ đồng.
}
\end{bt}

\begin{bt}%[9H9C29-3]
% nguon: Đề tham khảo Tuyển sinh 10 năm học 2026-2027 Sở GD&ĐT TP.HCM - Cụm 13 Đề 1 (Đề số 37)
Cho điểm $A$ nằm ngoài đường tròn $(O; R)$ sao cho $OA = 2R$. Kẻ hai tiếp tuyến $AB, AC$ với đường tròn $(O; R)$ ($B, C$ là các tiếp điểm), tia $AO$ cắt dây $BC$ tại $I$. Điểm $H$ thuộc đoạn thẳng $BI$ ($H$ khác $B$ và $I$). Đường thẳng $d$ vuông góc với $OH$ tại $H$; đường thẳng $d$ cắt $AB, AC$ lần lượt tại $P$ và $Q$.
\begin{enumerate}[a)]
	\item Chứng minh tứ giác $OHBP$ nội tiếp đường tròn.
	\item Chứng minh rằng $OP = OQ$.
	\item Khi $H$ là trung điểm của đoạn thẳng $BI$, tính độ dài đoạn thẳng $BC$ và diện tích của tam giác $OPQ$ theo $R$.
\end{enumerate}
\loigiai{
\begin{center}
\begin{tikzpicture}[>=stealth,scale=1.1]
	\def\R{2.2}
	\coordinate (O) at (0,0);
	\draw[thick] (O) circle (\R);
	\coordinate (A) at (2*\R,0);
	% Tiếp tuyến AB, AC: cos(AOB) = R/(2R) = 0.5 -> góc 60 độ
	\coordinate (B) at ({\R*cos(60)},{\R*sin(60)});
	\coordinate (C) at ({\R*cos(-60)},{\R*sin(-60)});
	\draw (A) -- (B) -- (O) -- (C) -- (A);
	\draw (O) -- (A);
	\draw (B) -- (C);
	\coordinate (I) at (intersection of B--C and O--A);
	% H trung điểm BI
	\coordinate (H) at ($(B)!0.5!(I)$);
	\draw (O) -- (H);
	% Đường thẳng d qua H vuông góc OH
	\coordinate (P) at ($(H)!2!90:(O)$);
	\coordinate (Q) at ($(H)!2!-90:(O)$);
	\coordinate (P_real) at (intersection of H--P and A--B);
	\coordinate (Q_real) at (intersection of H--Q and A--C);
	\draw (P_real) -- (Q_real);
	\draw (O) -- (P_real);
	\draw (O) -- (Q_real);
	% Ký hiệu
	\fill (O) circle (1.5pt) node[left] {$O$}
		  (A) circle (1.5pt) node[right] {$A$}
		  (B) circle (1.5pt) node[above] {$B$}
		  (C) circle (1.5pt) node[below] {$C$}
		  (I) circle (1.5pt) node[below right] {$I$}
		  (H) circle (1.5pt) node[above right] {$H$}
		  (P_real) circle (1.5pt) node[above right] {$P$}
		  (Q_real) circle (1.5pt) node[below right] {$Q$};
\end{tikzpicture}
\end{center}
\begin{enumerate}[a)]
	\item Vì $AB$ là tiếp tuyến của $(O)$ tại $B$ nên $AB \perp OB$, do đó $\widehat{OBP} = 90^\circ$.
	Tam giác $OBP$ vuông tại $B$ nên nội tiếp đường tròn đường kính $OP$.
	Lại có $d \perp OH$ tại $H$ nên $\widehat{OHP} = 90^\circ$.
	Tam giác $OHP$ vuông tại $H$ nên nội tiếp đường tròn đường kính $OP$.
	Suy ra bốn điểm $O, H, B, P$ cùng thuộc đường tròn đường kính $OP$, hay tứ giác $OHBP$ nội tiếp.
	\item Theo tính chất hai tiếp tuyến cắt nhau, ta có $AB = AC$ và $OB = OC = R$.
	Do đó $OA$ là đường trung trực của đoạn thẳng $BC$, suy ra $OA \perp BC$ tại $I$.
	Tương tự câu a:
	\begin{itemize}
		\item Vì $AC$ là tiếp tuyến của $(O)$ tại $C$ nên $\widehat{OCQ} = 90^\circ$, tam giác $OCQ$ nội tiếp đường tròn đường kính $OQ$.
		\item Vì $OH \perp HQ$ nên $\widehat{OHQ} = 90^\circ$, tam giác $OHQ$ nội tiếp đường tròn đường kính $OQ$.
	\end{itemize}
	Suy ra bốn điểm $O, H, C, Q$ cùng thuộc đường tròn đường kính $OQ$, hay tứ giác $OHCQ$ nội tiếp.
	Từ các tứ giác nội tiếp, ta có:
	\begin{itemize}
		\item Tứ giác $OHBP$ nội tiếp $\implies \widehat{OPH} = \widehat{OBH}$ (hai góc nội tiếp cùng chắn cung $OH$).
		\item Tứ giác $OHCQ$ nội tiếp $\implies \widehat{OQH} = \widehat{OCH}$ (hai góc nội tiếp cùng chắn cung $OH$).
	\end{itemize}
	Mà tam giác $OBC$ cân tại $O$ (do $OB = OC = R$) nên $\widehat{OBH} = \widehat{OCH}$.
	Suy ra $\widehat{OPH} = \widehat{OQH}$, hay $\widehat{OPQ} = \widehat{OQP}$.
	Do đó tam giác $OPQ$ cân tại $O$, suy ra $OP = OQ$.
	\item Trong tam giác vuông $OBA$ có $\cos \widehat{BOA} = \dfrac{OB}{OA} = \dfrac{R}{2R} = \dfrac{1}{2} \implies \widehat{BOA} = 60^\circ$.
	Xét tam giác vuông $OBI$ tại $I$:
	$$OI = OB \cdot \cos 60^\circ = \dfrac{R}{2}; \quad BI = OB \cdot \sin 60^\circ = \dfrac{R\sqrt{3}}{2}.$$
	Vì $OA$ là đường trung trực của $BC$ nên $I$ là trung điểm $BC$, suy ra:
	$$BC = 2BI = 2 \cdot \dfrac{R\sqrt{3}}{2} = R\sqrt{3}.$$
	Khi $H$ là trung điểm của $BI$, ta có:
	$$IH = \dfrac{BI}{2} = \dfrac{R\sqrt{3}}{4}.$$
	Áp dụng định lý Pythagore trong tam giác vuông $OIH$, ta có:
	$$OH^2 = OI^2 + IH^2 = \left(\dfrac{R}{2}\right)^2 + \left(\dfrac{R\sqrt{3}}{4}\right)^2 = \dfrac{R^2}{4} + \dfrac{3R^2}{16} = \dfrac{7R^2}{16} \implies OH = \dfrac{R\sqrt{7}}{4}.$$
	Xét $\triangle OHP$ và $\triangle OIB$ có:
	$$\widehat{OHP} = \widehat{OIB} = 90^\circ; \quad \widehat{OPH} = \widehat{OBI}\text{ (chứng minh câu b)}.$$
	Do đó $\triangle OHP \sim \triangle OIB$ (g-g).
	Suy ra:
	$$\dfrac{PH}{IB} = \dfrac{OH}{OI} \implies PH = \dfrac{OH \cdot IB}{OI} = \dfrac{\dfrac{R\sqrt{7}}{4} \cdot \dfrac{R\sqrt{3}}{2}}{\dfrac{R}{2}} = \dfrac{R\sqrt{21}}{4}.$$
	Vì tam giác $OPQ$ cân tại $O$ có $OH$ là đường cao nên $OH$ đồng thời là đường trung tuyến, suy ra:
	$$PQ = 2PH = 2 \cdot \dfrac{R\sqrt{21}}{4} = \dfrac{R\sqrt{21}}{2}.$$
	Diện tích tam giác $OPQ$ là:
	$$S_{\triangle OPQ} = \dfrac{1}{2} OH \cdot PQ = \dfrac{1}{2} \cdot \dfrac{R\sqrt{7}}{4} \cdot \dfrac{R\sqrt{21}}{2} = \dfrac{7R^2\sqrt{3}}{16}.$$
	Vậy $BC = R\sqrt{3}$ và $S_{\triangle OPQ} = \dfrac{7R^2\sqrt{3}}{16}$.
\end{enumerate}
}
\end{bt}

% ==============================================================================
% CỤM 13 - ĐỀ THAM KHẢO 2 (ĐỀ SỐ 38)
% ==============================================================================

\begin{bt}%[9D6H18-2]
% nguon: Đề tham khảo Tuyển sinh 10 năm học 2026-2027 Sở GD&ĐT TP.HCM - Cụm 13 Đề 2 (Đề số 38)
Cho hàm số $y = 0{,}5x^2$ có đồ thị là parabol $(P)$.
\begin{enumerate}[a)]
	\item Vẽ đồ thị $(P)$ trên mặt phẳng toạ độ $Oxy$.
	\item Tìm toạ độ điểm $M$ thuộc đồ thị $(P)$ có hoành độ bằng $4$.
\end{enumerate}
\loigiai{
\begin{enumerate}[a)]
	\item Bảng giá trị của hàm số:
	\begin{center}
		\begin{tabular}{|c|c|c|c|c|c|}
			\hline
			$x$ & $-4$ & $-2$ & $0$ & $2$ & $4$ \\ \hline
			$y = 0{,}5x^2$ & $8$ & $2$ & $0$ & $2$ & $8$ \\ \hline
		\end{tabular}
	\end{center}
	Đồ thị $(P)$ là một đường parabol đi qua các điểm $(-4; 8)$, $(-2; 2)$, $(0; 0)$, $(2; 2)$, $(4; 8)$, có đỉnh tại gốc toạ độ $O(0;0)$ và nhận trục tung $Oy$ làm trục đối xứng.
	\begin{center}
	\begin{tikzpicture}[>=stealth,scale=0.7]
		\draw[->] (-5,0) -- (5,0) node[below] {$x$};
		\draw[->] (0,-1) -- (0,9) node[left] {$y$};
		\node[below left] at (0,0) {$O$};
		\draw[dashed] (-4,0) node[below] {$-4$} -- (-4,8) -- (0,8) node[left] {$8$} -- (4,8) -- (4,0) node[below] {$4$};
		\draw[dashed] (-2,0) node[below] {$-2$} -- (-2,2) -- (0,2) node[left] {$2$} -- (2,2) -- (2,0) node[below] {$2$};
		\draw[thick,smooth,samples=100,domain=-4.2:4.2] plot(\x,{0.5*(\x)^2});
		\fill (-4,8) circle (1.5pt) (-2,2) circle (1.5pt) (0,0) circle (1.5pt) (2,2) circle (1.5pt) (4,8) circle (1.5pt);
		\node[right] at (2.5,5) {$(P)\colon y = 0{,}5x^2$};
	\end{tikzpicture}
	\end{center}
	\item Vì điểm $M$ thuộc $(P)$ và có hoành độ bằng $4$ nên thay $x = 4$ vào phương trình hàm số ta được:
	$$y = 0{,}5 \cdot 4^2 = 0{,}5 \cdot 16 = 8.$$
	Vậy toạ độ điểm cần tìm là $M(4; 8)$.
\end{enumerate}
}
\end{bt}

\begin{bt}%[9D6H20-2]
% nguon: Đề tham khảo Tuyển sinh 10 năm học 2026-2027 Sở GD&ĐT TP.HCM - Cụm 13 Đề 2 (Đề số 38)
Cho phương trình $-3x^2 + 2x + 4 = 0$.
\begin{enumerate}[a)]
	\item Chứng minh phương trình luôn có hai nghiệm phân biệt $x_1, x_2$.
	\item Không giải phương trình, hãy tính giá trị của biểu thức $A = \dfrac{2x_1}{x_2} + \dfrac{2x_2}{x_1}$.
\end{enumerate}
\loigiai{
\begin{enumerate}[a)]
	\item Phương trình $-3x^2 + 2x + 4 = 0$ có các hệ số $a = -3$, $b = 2$, $c = 4$.
	Biệt thức $\Delta$ của phương trình là:
	$$\Delta = b^2 - 4ac = 2^2 - 4 \cdot (-3) \cdot 4 = 4 + 48 = 52 > 0.$$
	(hoặc do $a \cdot c = (-3) \cdot 4 = -12 < 0$ nên phương trình có hai nghiệm trái dấu).
	Vì $\Delta > 0$ nên phương trình luôn có hai nghiệm phân biệt $x_1, x_2$.
	\item Theo hệ thức Viète, ta có:
	$$\begin{cases} x_1 + x_2 = -\dfrac{b}{a} = -\dfrac{2}{-3} = \dfrac{2}{3} \\ x_1 x_2 = \dfrac{c}{a} = \dfrac{4}{-3} = -\dfrac{4}{3}. \end{cases}$$
	Ta có:
	$$x_1^2 + x_2^2 = (x_1 + x_2)^2 - 2x_1 x_2 = \left(\dfrac{2}{3}\right)^2 - 2 \cdot \left(-\dfrac{4}{3}\right) = \dfrac{4}{9} + \dfrac{8}{3} = \dfrac{28}{9}.$$
	Khi đó biểu thức $A$ là:
	$$A = \dfrac{2x_1}{x_2} + \dfrac{2x_2}{x_1} = \dfrac{2(x_1^2 + x_2^2)}{x_1 x_2} = \dfrac{2 \cdot \dfrac{28}{9}}{-\dfrac{4}{3}} = \dfrac{\dfrac{56}{9}}{-\dfrac{4}{3}} = \dfrac{56}{9} \cdot \left(-\dfrac{3}{4}\right) = -\dfrac{14}{3}.$$
	Vậy $A = -\dfrac{14}{3}$.
\end{enumerate}
}
\end{bt}

\begin{bt}%[9D8H26-3]
% nguon: Đề tham khảo Tuyển sinh 10 năm học 2026-2027 Sở GD&ĐT TP.HCM - Cụm 13 Đề 2 (Đề số 38)
Lớp 9A có $2$ bạn nữ hát hay là Ngân và Hoa, $2$ bạn nam hát hay là Kiệt và Dũng. Cô giáo chủ nhiệm lớp muốn chọn ra $2$ bạn để hát song ca trong lễ bế giảng năm học.
\begin{enumerate}[a)]
	\item Hãy liệt kê các cách chọn ngẫu nhiên $2$ bạn để hát song ca.
	\item Tính xác suất của mỗi biến cố sau:
	\begin{itemize}
		\item $A$: ``Trong $2$ bạn được chọn có $1$ bạn nam và $1$ bạn nữ''.
		\item $B$: ``Trong $2$ bạn được chọn có bạn Kiệt''.
	\end{itemize}
\end{enumerate}
\loigiai{
\begin{enumerate}[a)]
	\item Ký hiệu các bạn: Ngân (N), Hoa (H), Kiệt (K), Dũng (D).
	Các kết quả có thể xảy ra khi chọn $2$ bạn là:
	$$\Omega = \{(N, H); (N, K); (N, D); (H, K); (H, D); (K, D)\}.$$
	Số phần tử của không gian mẫu là $n(\Omega) = 6$.
	\item Tính xác suất của từng biến cố:
	\begin{itemize}
		\item Xét biến cố $A$: ``Trong $2$ bạn được chọn có $1$ bạn nam và $1$ bạn nữ''.
		Các kết quả thuận lợi cho biến cố $A$ là: $(N, K); (N, D); (H, K); (H, D)$.
		Suy ra $n(A) = 4$.
		Xác suất của biến cố $A$ là:
		$$P(A) = \dfrac{n(A)}{n(\Omega)} = \dfrac{4}{6} = \dfrac{2}{3}.$$
		\item Xét biến cố $B$: ``Trong $2$ bạn được chọn có bạn Kiệt''.
		Các kết quả thuận lợi cho biến cố $B$ là: $(N, K); (H, K); (K, D)$.
		Suy ra $n(B) = 3$.
		Xác suất của biến cố $B$ là:
		$$P(B) = \dfrac{n(B)}{n(\Omega)} = \dfrac{3}{6} = \dfrac{1}{2}.$$
	\end{itemize}
\end{enumerate}
}
\end{bt}

\begin{bt}%[9D1H3-4]
% nguon: Đề tham khảo Tuyển sinh 10 năm học 2026-2027 Sở GD&ĐT TP.HCM - Cụm 13 Đề 2 (Đề số 38)
Theo Tổ chức Y tế Thế giới (WHO), dung dịch cồn $70^\circ$ được khuyến nghị đảm bảo tiêu diệt các loại virus, vi khuẩn gây hại. Để đảm bảo an toàn cho lớp học, cô Lan cùng nhóm học sinh đã pha $6$ lít cồn $70^\circ$ từ hai loại cồn $90^\circ$ và cồn $60^\circ$ để các bạn sát khuẩn tay. Hỏi cô Lan đã pha cồn $90^\circ$ và cồn $60^\circ$ theo tỉ lệ thể tích nào để được cồn $70^\circ$?
Biết: $\text{Độ cồn} = \dfrac{V_{\text{cồn nguyên chất}}}{V_{\text{dung dịch}}} \cdot 100^\circ$.
\loigiai{
Gọi $x$ (lít) là thể tích cồn $90^\circ$ và $y$ (lít) là thể tích cồn $60^\circ$ cần dùng ($0 < x, y < 6$).
Vì tổng thể tích dung dịch sau khi pha là $6$ lít nên ta có phương trình:
$$x + y = 6. \quad (1)$$
Lượng cồn nguyên chất trong $x$ lít cồn $90^\circ$ là:
$$\dfrac{90}{100}x = 0{,}9x\text{ (lít)}.$$
Lượng cồn nguyên chất trong $y$ lít cồn $60^\circ$ là:
$$\dfrac{60}{100}y = 0{,}6y\text{ (lít)}.$$
Lượng cồn nguyên chất trong $6$ lít cồn $70^\circ$ thu được là:
$$\dfrac{70}{100} \cdot 6 = 4{,}2\text{ (lít)}.$$
Vì tổng lượng cồn nguyên chất trước và sau khi pha không đổi nên ta có phương trình:
$$0{,}9x + 0{,}6y = 4{,}2. \quad (2)$$
Từ $(1)$ và $(2)$ ta có hệ phương trình:
$$\begin{cases} x + y = 6 \\ 0{,}9x + 0{,}6y = 4{,}2 \end{cases} \iff \begin{cases} 0{,}6x + 0{,}6y = 3{,}6 \\ 0{,}9x + 0{,}6y = 4{,}2 \end{cases} \iff \begin{cases} 0{,}3x = 0{,}6 \\ x + y = 6 \end{cases} \iff \begin{cases} x = 2 \\ y = 4 \end{cases}\text{ (thoả mãn)}.$$
Tỉ lệ thể tích cồn $90^\circ$ và cồn $60^\circ$ cần pha là:
$$\dfrac{x}{y} = \dfrac{2}{4} = \dfrac{1}{2}.$$
Vậy cô Lan đã pha cồn $90^\circ$ và cồn $60^\circ$ theo tỉ lệ thể tích là $1 : 2$.
}
\end{bt}

\begin{bt}%[9H0V31-3]
% nguon: Đề tham khảo Tuyển sinh 10 năm học 2026-2027 Sở GD&ĐT TP.HCM - Cụm 13 Đề 2 (Đề số 38)
Để tổ chức sinh nhật cho con gái, chị Thanh đã đặt thợ làm bánh tại cửa hàng Bakery làm một chiếc bánh hai tầng hình trụ, mỗi tầng cao $15\text{ cm}$, bán kính tầng trên là $15\text{ cm}$, đường kính tầng dưới là $40\text{ cm}$.
\begin{enumerate}[a)]
	\item Tính thể tích chiếc bánh. (Làm tròn kết quả đến hai chữ số thập phân).
	\item Khi chiếc bánh hoàn thành, người thợ có tất cả bao nhiêu diện tích bề mặt bên ngoài để trang trí bánh? (Bao gồm diện tích xung quanh hai tầng, mặt trên của tầng nhỏ và phần vành khăn hở ra của tầng lớn, đáy bánh đặt trên khay không trang trí. Làm tròn kết quả đến hai chữ số thập phân).
\end{enumerate}
\loigiai{
\begin{enumerate}[a)]
	\item Tầng bánh dưới hình trụ có bán kính đáy là:
	$$r_2 = 40 : 2 = 20\text{ (cm)}.$$
	Thể tích của tầng bánh trên là:
	$$V_1 = \pi r_1^2 h = \pi \cdot 15^2 \cdot 15 = 3375\pi\text{ (cm}^3\text{)}.$$
	Thể tích của tầng bánh dưới là:
	$$V_2 = \pi r_2^2 h = \pi \cdot 20^2 \cdot 15 = 6000\pi\text{ (cm}^3\text{)}.$$
	Tổng thể tích của chiếc bánh là:
	$$V = V_1 + V_2 = 3375\pi + 6000\pi = 9375\pi \approx 29452{,}43\text{ (cm}^3\text{)}.$$
	Vậy thể tích của chiếc bánh xấp xỉ $29452{,}43\text{ cm}^3$.
	\item Diện tích xung quanh của tầng bánh trên là:
	$$S_{xq1} = 2\pi r_1 h = 2\pi \cdot 15 \cdot 15 = 450\pi\text{ (cm}^2\text{)}.$$
	Diện tích xung quanh của tầng bánh dưới là:
	$$S_{xq2} = 2\pi r_2 h = 2\pi \cdot 20 \cdot 15 = 600\pi\text{ (cm}^2\text{)}.$$
	Diện tích bề mặt nằm ngang cần trang trí gồm mặt trên của tầng nhỏ và phần vành khăn của tầng lớn, tổng diện tích này đúng bằng diện tích hình tròn đáy của tầng dưới:
	$$S_{\text{ngang}} = \pi r_2^2 = \pi \cdot 20^2 = 400\pi\text{ (cm}^2\text{)}.$$
	Tổng diện tích bề mặt cần trang trí là:
	$$S = S_{xq1} + S_{xq2} + S_{\text{ngang}} = 450\pi + 600\pi + 400\pi = 1450\pi \approx 4555{,}31\text{ (cm}^2\text{)}.$$
	Vậy diện tích bề mặt để trang trí bánh là khoảng $4555{,}31\text{ cm}^2$.
\end{enumerate}
}
\end{bt}

\begin{bt}%[9D8H26-3]
% nguon: Đề tham khảo Tuyển sinh 10 năm học 2026-2027 Sở GD&ĐT TP.HCM - Cụm 13 Đề 2 (Đề số 38)
Biểu đồ hình quạt tròn biểu diễn điểm số của học sinh lớp 9A qua đợt kiểm tra thường xuyên môn Toán như sau: điểm $6$ chiếm $15\%$, điểm $7$ chiếm $20\%$, điểm $8$ chiếm $25\%$, điểm $9$ chiếm $30\%$ và điểm $10$ chiếm $10\%$.
\begin{enumerate}[a)]
	\item Biết sĩ số của lớp 9A là $40$ học sinh. Tính tổng số học sinh đạt điểm $8$ và điểm $9$.
	\item Chọn ngẫu nhiên một bạn trong lớp tham gia cuộc thi giao lưu. Tính xác suất của biến cố: ``Bạn được chọn có điểm kiểm tra Toán thấp nhất là $8$ điểm''.
\end{enumerate}
\loigiai{
\begin{enumerate}[a)]
	\item Số học sinh đạt điểm $8$ là:
	$$40 \cdot 25\% = 10\text{ (học sinh)}.$$
	Số học sinh đạt điểm $9$ là:
	$$40 \cdot 30\% = 12\text{ (học sinh)}.$$
	Tổng số học sinh đạt điểm $8$ và điểm $9$ là:
	$$10 + 12 = 22\text{ (học sinh)}.$$
	\item Chọn ngẫu nhiên $1$ bạn trong lớp có $40$ học sinh nên số phần tử của không gian mẫu là $n(\Omega) = 40$.
	Biến cố ``Bạn được chọn có điểm kiểm tra Toán thấp nhất là $8$ điểm'' có nghĩa là bạn đó đạt điểm $8$, điểm $9$ hoặc điểm $10$.
	Số học sinh đạt điểm $10$ là:
	$$40 \cdot 10\% = 4\text{ (học sinh)}.$$
	Số kết quả thuận lợi cho biến cố là:
	$$m = 10 + 12 + 4 = 26.$$
	Xác suất của biến cố là:
	$$P = \dfrac{m}{n(\Omega)} = \dfrac{26}{40} = \dfrac{13}{20}.$$
\end{enumerate}
}
\end{bt}

\begin{bt}%[9H9C29-3]
% nguon: Đề tham khảo Tuyển sinh 10 năm học 2026-2027 Sở GD&ĐT TP.HCM - Cụm 13 Đề 2 (Đề số 38)
Cho tam giác $ABC$ nhọn ($AB < AC$) nội tiếp đường tròn $(O; R)$, hai đường cao $BE$ và $AF$ cắt nhau tại trực tâm $H$. Vẽ đường kính $AD$ của đường tròn $(O)$. Qua $H$ vẽ đường thẳng $d$ vuông góc với $AD$ tại $K$; đường thẳng $d$ cắt các cạnh $AB, AC$ và đường thẳng $BC$ lần lượt tại $M, N$ và $S$.
\begin{enumerate}[a)]
	\item Chứng minh $4$ điểm $A, E, H, K$ cùng thuộc một đường tròn. Xác định tâm $I$ của đường tròn này.
	\item Chứng minh $\triangle AKM \sim \triangle ABD$ và $SM \cdot SN = SB \cdot SC$.
	\item Chứng minh $SI \perp OI$.
\end{enumerate}
\loigiai{
\begin{center}
\begin{tikzpicture}[>=stealth,scale=1.1]
	\def\R{2.6}
	\coordinate (O) at (0,0);
	\draw[thick] (O) circle (\R);
	\coordinate (A) at (-0.6,2.53);
	\coordinate (B) at (-2.4,-0.99);
	\coordinate (C) at (2.2,-1.38);
	\draw (A) -- (B) -- (C) -- cycle;
	% Đường cao AF, BE
	\coordinate (F) at ($(B)!(A)!(C)$);
	\coordinate (E) at ($(A)!(B)!(C)$);
	\draw (A) -- (F);
	\draw (B) -- (E);
	\coordinate (H) at (intersection of A--F and B--E);
	% Đường kính AD
	\coordinate (D) at ($(O)!-1!(A)$);
	\draw (A) -- (D);
	\draw (B) -- (D);
	\draw (C) -- (D);
	% Đường thẳng d qua H vuông góc AD tại K
	\coordinate (K) at ($(A)!(H)!(D)$);
	\coordinate (M) at (intersection of H--K and A--B);
	\coordinate (N) at (intersection of H--K and A--C);
	\coordinate (S) at (intersection of H--K and B--C);
	\draw (S) -- (N);
	\draw (S) -- (B);
	% I là trung điểm AH
	\coordinate (I) at ($(A)!0.5!(H)$);
	\draw (S) -- (I);
	\draw (O) -- (I);
	% Ký hiệu
	\fill (O) circle (1.5pt) node[below left] {$O$}
		  (A) circle (1.5pt) node[above] {$A$}
		  (B) circle (1.5pt) node[left] {$B$}
		  (C) circle (1.5pt) node[right] {$C$}
		  (H) circle (1.5pt) node[above right] {$H$}
		  (E) circle (1.5pt) node[above right] {$E$}
		  (F) circle (1.5pt) node[below] {$F$}
		  (D) circle (1.5pt) node[below] {$D$}
		  (K) circle (1.5pt) node[above right] {$K$}
		  (M) circle (1.5pt) node[above left] {$M$}
		  (N) circle (1.5pt) node[above right] {$N$}
		  (S) circle (1.5pt) node[below left] {$S$}
		  (I) circle (1.5pt) node[left] {$I$};
\end{tikzpicture}
\end{center}
\begin{enumerate}[a)]
	\item Vì $BE \perp AC$ tại $E$ nên $\widehat{AEH} = 90^\circ$, suy ra tam giác $AEH$ vuông tại $E$ nên nội tiếp đường tròn đường kính $AH$.
	Vì $d \perp AD$ tại $K$ nên $\widehat{AKH} = 90^\circ$, suy ra tam giác $AKH$ vuông tại $K$ nên nội tiếp đường tròn đường kính $AH$.
	Do đó bốn điểm $A, E, H, K$ cùng thuộc đường tròn đường kính $AH$.
	Tâm $I$ của đường tròn này là trung điểm của đoạn thẳng $AH$.
	\item Ta có $\widehat{ABD} = 90^\circ$ (góc nội tiếp chắn nửa đường tròn đường kính $AD$).
	Xét $\triangle AKM$ và $\triangle ABD$ có:
	\begin{itemize}
		\item $\widehat{BAD}$ là góc chung;
		\item $\widehat{AKM} = \widehat{ABD} = 90^\circ$.
	\end{itemize}
	Do đó $\triangle AKM \sim \triangle ABD$ (g-g).
	Suy ra $\widehat{AMK} = \widehat{ADB}$.
	Mặt khác, trong đường tròn $(O)$, $\widehat{ADB} = \widehat{ACB}$ (hai góc nội tiếp cùng chắn cung $AB$).
	Do đó $\widehat{AMK} = \widehat{ACB}$.
	Mà $\widehat{SMB} = \widehat{AMK}$ (hai góc đối đỉnh), suy ra $\widehat{SMB} = \widehat{SCN}$.
	Xét $\triangle SMB$ và $\triangle SNC$ có:
	\begin{itemize}
		\item $\widehat{CSN}$ là góc chung;
		\item $\widehat{SMB} = \widehat{SCN}$ (chứng minh trên).
	\end{itemize}
	Do đó $\triangle SMB \sim \triangle SNC$ (g-g).
	Suy ra $\dfrac{SM}{SC} = \dfrac{SB}{SN} \implies SM \cdot SN = SB \cdot SC$.
	\item Ta có:
	\begin{itemize}
		\item $\widehat{ACD} = 90^\circ$ (góc nội tiếp chắn nửa đường tròn $(O)$) $\implies CD \perp AC$.
		Lại có $BH \perp AC$ (do $BE$ là đường cao) $\implies BH \parallel CD$.
		\item $\widehat{ABD} = 90^\circ$ $\implies BD \perp AB$.
		Lại có $CH \perp AB$ (do $H$ là trực tâm) $\implies CH \parallel BD$.
	\end{itemize}
	Tứ giác $BHCD$ có $BH \parallel CD$ và $CH \parallel BD$ nên là hình bình hành.
	Gọi $Q$ là giao điểm của hai đường chéo $BC$ và $HD$, suy ra $Q$ là trung điểm của $HD$ và $Q \in BC$.
	Trong tam giác $AHD$:
	\begin{itemize}
		\item $I$ là trung điểm của $AH$;
		\item $Q$ là trung điểm của $HD$.
	\end{itemize}
	Suy ra $IQ$ là đường trung bình của tam giác $AHD$, do đó $IQ \parallel AD$.
	Mà $d \perp AD$ tại $K$ nên $IQ \perp d$, hay $IQ \perp SH$.
	Trong tam giác $SIQ$:
	\begin{itemize}
		\item $SH \perp IQ$ (chứng minh trên), nên $SH$ là đường cao kẻ từ đỉnh $S$;
		\item $AI \perp SQ$ (vì $AH \perp BC$ tại $F$, mà $I \in AH$ và $S, Q \in BC$), nên $AI$ là đường cao kẻ từ đỉnh $I$.
	\end{itemize}
	Hai đường cao $SH$ và $AI$ cắt nhau tại $H$, suy ra $H$ là trực tâm của tam giác $SIQ$.
	Do đó $QH \perp SI$. Mà $Q, H, D$ thẳng hàng nên $HD \perp SI$.
	Mặt khác, trong tam giác $AHD$, $I$ là trung điểm $AH$ và $O$ là trung điểm $AD$ nên $IO$ là đường trung bình của tam giác $AHD$, suy ra $IO \parallel HD$.
	Vì $IO \parallel HD$ và $HD \perp SI$ nên $SI \perp OI$.
\end{enumerate}
}
\end{bt}
